\chapter{Reproducibilidad}
\label{anexo:reproducibilidad}

\section{Identidad de los estudios}

La evidencia empírica de este trabajo procede de cuatro estudios encadenados: tres \textit{Model
Studies} en los que el ganador de cada pasada es el \textit{baseline} de la siguiente, y un
\textit{Portfolio Study} que optimiza la cartera del último sin reentrenar nada. Sus identificadores
permiten localizar cualquier cifra citada en el texto.

\begin{description}[leftmargin=3.2cm,style=nextline,font=\normalfont\ttfamily]
  \item[Model Study 1] \studyidA{} · ganador \artefacto{run-6eaa47a0597b} · catálogo v6
  \item[Model Study 2] \studyidB{} · ganador \artefacto{run-2dc586be8653} · catálogo v6
  \item[Model Study 3] \studyidC{} · ganador \artefacto{run-f134d7eb9e06} · catálogo v7
  \item[Portfolio Study] \studyportfolio{} · 1.728 carteras evaluadas
\end{description}

El Model Study 3 es el estudio de referencia del documento: de él sale toda la evidencia predictiva.
Sus identificadores de datos son los siguientes.

\begin{description}[leftmargin=3.2cm,style=nextline,font=\normalfont\ttfamily]
  \item[dataset\_hash] \artefacto{b9134b218e3bf7fc\ldots{}a69df0a92653fbd3}
  \item[evaluation\_key] \artefacto{0088d08579c1e65f\ldots{}5dbf2fa71de0a27f}
  \item[catalog\_hash] \artefacto{845d55249866f2f0\ldots{}049f59a8b1a92e7}
  \item[catalog\_version] 7
\end{description}

Las tres pasadas emplearon la misma ventana de selección --- 117 cohortes mensuales entre 2015 y
2024 --- y reservaron 2025--2026 sin que interviniera en ninguna decisión. La métrica de selección
fue \texttt{rank\_ic\_only} en los tres casos y \texttt{information\_ratio} en el Portfolio Study,
cuya rejilla se calculó además sobre una serie recortada en 2024 para que ninguna de las 1.728
combinaciones pudiera observar la era reservada.

Las pasadas 1 y 2 corrieron bajo la versión 6 del catálogo y la tercera bajo la 7, de modo que no
son estrictamente comparables entre sí en sus reglas de desempate. El alcance de esa discrepancia se
analiza en la Sección~\ref{sec:reproducibilidad}.

\section{Regeneración de figuras y tablas}

Las figuras y tablas del manuscrito no se dibujan a mano: se generan desde los artefactos
persistidos mediante un único comando, ejecutado desde la raíz del repositorio.

\begin{quote}
\ttfamily
python latex/scripts/export\_study\_assets.py \\
\hspace*{2em}--study-id study-20260814-095144-5ec17b78 \\
\hspace*{2em}--chain-study-id study-20260812-163136-1b104667 \\
\hspace*{2em}--chain-study-id study-20260813-103456-aa733655 \\
\hspace*{2em}--chain-study-id study-20260814-095144-5ec17b78 \\
\hspace*{2em}--portfolio-study-id study-20260814-135754-fdbdf2c5
\end{quote}

Los tres \texttt{--chain-study-id} se declaran en orden de ejecución y el último coincide con
\texttt{--study-id}: de ellos salen las figuras y tablas de la cadena. La separación entre
\texttt{--study-id} y \texttt{--portfolio-study-id} es deliberada y hace explícita la procedencia de
cada activo: del primero sale todo lo \emph{predictivo} --- Rank-IC, agentes, meta-agente,
decisiones, robustez y atribución --- y del segundo todo lo \emph{económico} --- rejilla de
carteras, patrimonio, órdenes, métricas anuales y perfiles ---, leído de la evidencia del ganador
sobre la serie completa.

El script escribe los PNG a 300\,ppp y los cuerpos de tabla en \artefacto{latex/assets/}, junto al
resto de capítulos, y actualiza \artefacto{latex/asset\_manifest.json} con la lista de
artefactos consumidos, separando las fuentes predictivas de las económicas. Ese manifiesto es la
tabla de correspondencia entre cada salida del manuscrito y su origen: ninguna cifra generada por
script entra en el documento sin quedar registrada allí.

Dos elementos son la excepción y conviene declararlo: las Tablas~\ref{tab:defectos}
y~\ref{tab:versiones} del Anexo~\ref{anexo:desarrollo} se transcriben de la bitácora del proyecto
y no de un artefacto del estudio, porque documentan el proceso de desarrollo y no una medición.

Hay una segunda salvedad, sobre qué se puede regenerar y qué no. El repositorio versiona los
artefactos en formato JSON --- \artefacto{winner.json}, \artefacto{decisions.json},
\artefacto{robustness.json}, \artefacto{attribution.json}, \artefacto{portfolio\_winner.json} y los
\artefacto{summary.json} de cada evidencia ---, que son los que sostienen la mayoría de las cifras
citadas en el texto y permiten verificarlas directamente. Los \artefacto{.parquet}, en cambio,
quedan excluidos por tamaño, de modo que las tablas y figuras que se construyen a partir de ellos
--- cobertura, agentes, atribución local, colas, órdenes, métricas anuales, rejilla de carteras y
perfiles --- se conservan ya generadas en \artefacto{latex/assets/} pero no pueden recalcularse
desde un clon del repositorio sin volver a ejecutar el estudio. Se declara porque el comando
anterior sólo es reproducible en su totalidad sobre la instalación donde corrieron los estudios.

Un segundo script comprueba que el proyecto sea portable antes de subirlo:

\begin{quote}
\ttfamily
python latex/scripts/verify\_latex\_assets.py
\end{quote}

Verifica que no haya \textit{mojibake}, que todas las rutas de figura sean relativas y existan, que
cada tabla incluida exista, que toda referencia cruzada tenga destino, que no queden figuras ni
tablas generadas sin insertar y que no reaparezca ningún comando de bibliografía.

